% Chapter 4: Prompt Engineering That Scales
\chapter{Prompt Engineering That Scales}
\label{ch:prompting}

Prompt engineering for Claude Code differs from chatbot prompting
in a fundamental way: you are not crafting one-off queries but building
\emph{reusable systems} that run reliably across hundreds of sessions.
The techniques in this chapter optimize for repeatability, not cleverness.

\section{Three Highest-Leverage Techniques}
\label{sec:three-techniques}

\marginnote[Official]{Anthropic's prompt engineering guide identifies these as the techniques with the largest impact on output quality.\sourceurl{https://platform.claude.com/docs/en/build-with-claude/prompt-engineering/overview}}

\official{Anthropic's research consistently identifies three techniques that produce
the largest quality improvements:}

\begin{enumerate}
  \item \textbf{Tell Claude to think step by step.}
    Not just ``think carefully''---give explicit reasoning structure.
    ``First, identify the root cause. Then, list possible fixes.
    Finally, recommend the fix with the best tradeoff of effort vs.\ impact.''

  \item \textbf{Provide concrete examples (3--5 diverse).}
    Examples anchor Claude's understanding better than abstract descriptions.
    Include edge cases in your examples to demonstrate boundary handling.

  \item \textbf{Use structured output formats.}
    Specify the exact format you want: JSON schema, table structure,
    section headers. Claude follows explicit format instructions precisely.
\end{enumerate}

\subsection{Additional High-Value Techniques}

\begin{description}
  \item[XML tags] For multi-component prompts, use \code{<context>}, \code{<task>},
    \code{<constraints>} to separate concerns.
  \item[Rich inputs] Use \code{@file} to include files, paste screenshots,
    pipe data from \code{stdin}. Real data beats verbal descriptions.
  \item[Verification criteria] ``The single highest-leverage thing you can do
    is give Claude a way to verify its own work.''\sourceurl{https://code.claude.com/docs/en/best-practices}
\end{description}

\margintip{Verification criteria in CLAUDE.md: ``Always run tests after code changes. Never commit without passing CI.'' This one line improves every session.}

% -------------------------------------------------------------------
\section{Long Context Strategy}
\label{sec:long-context}

\marginnote[Official]{Claude Opus 4.6 supports 1M tokens (beta). Place documents at top, queries at bottom.\sourceurl{https://platform.claude.com/docs/en/build-with-claude/prompt-engineering/claude-prompting-best-practices}}

\official{For multi-document inputs, structure matters:}

\begin{enumerate}
  \item Place all reference documents at the \textbf{top} of the prompt.
  \item Place your question or task at the \textbf{bottom}.
  \item Ask Claude to \textbf{quote relevant passages first} before answering.
    This grounding technique can improve accuracy by 30\% on complex multi-document reasoning.
\end{enumerate}

\official{System prompts set the role; user prompts give the task.}
Claude follows user-turn instructions more reliably than system prompt instructions.
Use the system prompt for persona and constraints; use the user turn for
specific task instructions.

% -------------------------------------------------------------------
\section{Building a Precision Vocabulary}
\label{sec:precision}

\practitioner{Track how your requests could be stated more precisely.
Over time, develop a shared vocabulary between you and Claude.}

The pattern:
\begin{enumerate}
  \item After each interaction, note how the request could have been clearer.
  \item Identify recurring request patterns (create, diagnose, refactor, validate).
  \item Document these patterns in a project nomenclature.
  \item Use the precise version in future requests.
\end{enumerate}

\begin{fullwidth}
\begin{tabular}{@{}p{5cm}p{7cm}@{}}
\toprule
\textbf{Natural Language} & \textbf{Precise Version} \\
\midrule
``Fix the broken tests'' & \code{diagnose test failures in tests/auth/ \&\& fix root cause} \\
``Make this code better'' & \code{refactor process\_payment() to <20 lines, extract validation} \\
``Check if this works'' & \code{validate: run pytest, check coverage >80\%, verify edge cases} \\
\bottomrule
\end{tabular}
\end{fullwidth}

\margintip{Precision compounds. A team that develops shared vocabulary communicates 3x more efficiently with Claude after one month.}

% -------------------------------------------------------------------
\section{Extended Thinking}
\label{sec:thinking}

\marginnote[Official]{On by default. Toggle: Alt+T. Opus 4.6: adaptive thinking with effort controls.\sourceurl{https://platform.claude.com/docs/en/build-with-claude/extended-thinking}}

\official{Extended thinking lets Claude reason through complex problems
before responding.} On Opus 4.6, thinking is \textbf{adaptive}---the model
decides how much to think based on problem complexity.

\subsection{Effort Controls}

\begin{description}
  \item[\code{low}] Minimal thinking. Best for: simple lookups, formatting, boilerplate.
  \item[\code{medium}] Moderate analysis. Best for: standard coding tasks, reviews.
  \item[\code{high}] Deep reasoning. Best for: architecture decisions, complex debugging.
  \item[\code{max}] Maximum effort. Best for: strategic decisions, novel problems.
\end{description}

\official{Interleaved thinking}: Claude can think between tool calls,
enabling nuanced multi-step workflows where each step's reasoning
informs the next.

\subsection{When to Enable vs.\ Disable}

\begin{itemize}
  \item \textbf{Enable} for: complex multi-step reasoning, strategic decisions,
    code review where you want to see the thought process.
  \item \textbf{Disable} for: simple tasks, real-time applications,
    when latency matters more than depth.
\end{itemize}

% -------------------------------------------------------------------
\section{Cost Optimization}
\label{sec:cost}

\margincost{Prompt caching: 90\% savings on cached reads. Batch API: 50\% discount. Combined: 70\%+ savings.\sourceurl{https://platform.claude.com/docs/en/build-with-claude/prompt-caching}}

\subsection{Prompt Caching}

\official{Prompt caching reduces input costs by 90\% for repeated content.}

\begin{description}
  \item[Ephemeral (5-minute TTL)] Automatic for any content appearing in
    multiple requests within a short window.
  \item[Persistent (1-hour TTL)] Explicit cache breakpoints for content you know
    will be reused. Best for: tool definitions, system prompts, reference documents.
\end{description}

Cacheable content: tool definitions, system messages, examples, documentation,
images, PDFs.

\subsection{Batch API}

\official{50\% discount for asynchronous processing.} Submit up to 10,000
requests per batch. Results returned within 24 hours.
Combinable with caching for 70\%+ total savings.

Best for: bulk migrations, code review across many files, documentation generation.

\subsection{Model Selection}

\begin{fullwidth}
\begin{tabular}{@{}llp{6cm}@{}}
\toprule
\textbf{Model Class} & \textbf{Strength} & \textbf{Best For} \\
\midrule
Haiku & Speed, lowest cost & Simple tasks, high volume, formatting \\
Sonnet & Balance of speed and capability & Most coding tasks, daily development \\
Opus & Frontier reasoning & Complex architecture, multi-file refactoring \\
\bottomrule
\end{tabular}
\end{fullwidth}

\margincost{Pricing changes frequently. Check \url{https://platform.claude.com/docs/en/about-claude/pricing} for current rates.}

\practitioner{Default to Sonnet for development work. Escalate to Opus for
architecture decisions, complex debugging, and multi-file refactoring.
Use Haiku for bulk operations and simple formatting tasks.}

% -------------------------------------------------------------------
\section{Visual: Cost Optimization Decision Tree}

\begin{figure}[H]
\centering
\begin{tikzpicture}[
  decision/.style={draw=PrimaryBlue, fill=PrimaryBlue!8, diamond, aspect=2.2,
    font=\scriptsize, align=center, inner sep=1pt},
  action/.style={draw=TipGreen!80!black, fill=TipGreen!8, rounded corners=3pt,
    font=\scriptsize\bfseries, align=center, minimum width=1.8cm, minimum height=0.7cm},
  arr/.style={->, thick, DarkText},
  lbl/.style={font=\tiny, text=DarkText},
  >=Stealth
]
  % Root decision
  \node[decision] (q1) at (0,0) {Bulk operation?\\(>10 files)};

  % Yes branch — batch
  \node[decision] (q2) at (-3.5,-1.8) {Repeated\\content?};
  \node[action] (batch) at (-3.5,-3.5) {Batch API\\50\% off};
  \node[action] (batchcache) at (-5.5,-3.5) {Batch +\\Cache\\70\%+ off};

  % No branch — interactive
  \node[decision] (q3) at (3.5,-1.8) {Complex\\reasoning?};
  \node[action] (opus) at (5.5,-3.5) {Opus 4.6};
  \node[decision] (q4) at (1.5,-3.5) {Simple\\task?};
  \node[action] (haiku) at (0,-5.2) {Haiku 4.5};
  \node[action] (sonnet) at (3,-5.2) {Sonnet 4.6};

  % Arrows
  \draw[arr] (q1) -- (q2) node[midway, above, lbl] {Yes};
  \draw[arr] (q1) -- (q3) node[midway, above, lbl] {No};
  \draw[arr] (q2) -- (batchcache) node[midway, above, lbl] {Yes};
  \draw[arr] (q2) -- (batch) node[midway, right, lbl] {No};
  \draw[arr] (q3) -- (opus) node[midway, above, lbl] {Yes};
  \draw[arr] (q3) -- (q4) node[midway, right, lbl] {No};
  \draw[arr] (q4) -- (haiku) node[midway, above, lbl] {Yes};
  \draw[arr] (q4) -- (sonnet) node[midway, above, lbl] {No};
\end{tikzpicture}
\caption{Model and API selection decision tree for cost optimization.}
\label{fig:cost-decision-tree}
\end{figure}

\begin{keyinsight}
The most cost-effective pattern is not choosing the cheapest model---it is
structuring your workflow to maximize cache hits. A well-organized CLAUDE.md
with stable tool definitions can reduce per-session costs by 60--80\%
through caching alone.
\end{keyinsight}
